\begin{figure}[H]
    \centering
    \begin{tikzpicture}[font=\footnotesize, >=Latex, node distance=4mm,
        cstep/.style={draw, rounded corners, align=center, text width=0.72\linewidth,
                     minimum height=0.8cm, fill=gray!8}]
        \node[cstep] (px)   {\textbf{Pixel-scale calibration} (base objective (e.g. 4x), calibration slide)\\
                             calculates $\mu$m/pixel};
        \node[cstep, below=of px]   (derive) {\textbf{Derive other objectives}\\
                             scale by magnification ($M$) and numerical aperture ($NA$)};
        \node[cstep, below=of derive] (stage) {\textbf{Stage-scale calibration} (X and Y separately)\\
                             calculates $\mu$m/step per axis};
        \node[cstep, below=of stage] (back) {\textbf{Backlash calibration} (per axis, image comparison)\\
                              calculates backlash steps per axis};
        \node[cstep, below=of back] (focus) {\textbf{Depth-of-field (focus stacking)}\\
                             step spacing, search range};
        \draw[-{Latex}] (px)--(derive); \draw[-{Latex}] (derive)--(stage);
        \draw[-{Latex}] (stage)--(back); \draw[-{Latex}] (back)--(focus);
    \end{tikzpicture}
    \caption[Calibration workflow]{Calibration workflow. A base objective is calibrated against a slide, remaining objectives
    are derived from magnification and numerical aperture, and the stage scale, backlash, focus
    parameters, and soft limits are measured and stored in the configuration for reuse by the
    automated workflows.}
    \label{fig:calibration-flow}
\end{figure}
